% !TEX root = ../discretemath.tex

\section{Теорема Геринга/Менгера}

\subsection{Разделяющее множество}

\begin{Definition}{Разделяющее множество}{}
  В графе $G$ множество вершин $R \subset V$ называется
  \emph{разделяющим} для множеств $X \subset V$ и $Y \subset V$
  (которые могут пересекаться), если в подграфе $G[V \setminus R]$
  не существует пути между вершинами из $X \setminus R$
  и вершинами из $Y \setminus R$.
\end{Definition}

\begin{figure}[H]
\centering
\begin{tikzpicture}[
  >=Stealth, font=\small,
  vX/.style={circle, draw=blue!70!black, fill=blue!15, minimum size=8mm, inner sep=0pt},
  vR/.style={circle, draw=green!60!black, fill=green!20, minimum size=8mm, inner sep=0pt},
  vY/.style={circle, draw=red!70!black,  fill=red!15,  minimum size=8mm, inner sep=0pt},
  vN/.style={circle, draw=gray!60, fill=gray!10, minimum size=7mm, inner sep=0pt}]

\node[vX] (x1) at (0,  1.2) {$x_1$};
\node[vX] (x2) at (0, -1.2) {$x_2$};
\node[vR] (r1) at (3.2,  1.2) {$r_1$};
\node[vR] (r2) at (3.2, -1.2) {$r_2$};
\node[vY] (y1) at (6.4,  1.2) {$y_1$};
\node[vY] (y2) at (6.4, -1.2) {$y_2$};
\node[vN] (v)  at (1.6, 0) {};

\draw (x1)--(v)--(r1) (x2)--(v)--(r2);
\draw[gray!50] (x1) to[bend left=18]  (r1);
\draw[gray!50] (x2) to[bend right=18] (r2);
\draw (r1)--(y1) (r2)--(y2) (r1)--(y2) (r1)--(r2) (y1)--(y2);

\begin{scope}[on background layer]
  \draw[blue!40,  dashed, thick, rounded corners=6pt] (-0.6,-1.9) rectangle (0.6, 1.9);
  \draw[green!55!black, dashed, thick, rounded corners=6pt] (2.6,-1.9) rectangle (3.8, 1.9);
  \draw[red!50,   dashed, thick, rounded corners=6pt] (5.8,-1.9) rectangle (7.0, 1.9);
\end{scope}

\node[blue!70!black,  below] at (0,   -2.1) {$X$};
\node[green!60!black, below] at (3.2, -2.1) {$R$};
\node[red!70!black,   below] at (6.4, -2.1) {$Y$};
\end{tikzpicture}
\caption{Разделяющее множество $R$ (зелёное) отделяет $X$ (синее) от $Y$ (красное):
  в $G[V\setminus R]$ не существует пути из $X\setminus R$ в $Y\setminus R$.}
\end{figure}

\subsection{Теорема Геринга}

\begin{Theorem}{Геринга}{}
  Пусть задан граф $G$ и два множества его вершин $X, Y$.
  Если каждое $XY$-разделяющее множество содержит хотя бы $k$ вершин,
  то существует не менее $k$ попарно непересекающихся по вершинам
  $XY$-путей.
\end{Theorem}

\begin{Statement}{Эквивалентная переформулировка}{}
  Максимальное число попарно непересекающихся по вершинам $XY$-путей
  равно минимальному размеру $XY$-разделяющего множества.
\end{Statement}

\begin{proof}
  Индукция по $|V|$.

  \subsubsection*{База: $|V| = k$}

  $R = X$ — $XY$-разделяющее, поэтому $|X| \ge k$.
  Аналогично $|Y| \ge k$. При $|V|=k$ получаем $X=Y=V$,
  и из каждой вершины в себя есть путь длины $0$ — ровно $k$ непересекающихся путей. $\checkmark$

  \subsubsection*{Переход: $|V| > k$}

  \noindent\textbf{Случай 1.}
  $\exists\, R$ — $XY$-разделяющее, $|R|=k$,
  $X\setminus R \neq \emptyset$, $Y\setminus R \neq \emptyset$.

  Положим $V_X$ — вершины, достижимые из $X\setminus R$ в $G\setminus R$,
  $V_Y$ — из $Y\setminus R$; $G_X=G[V_X\cup R]$, $G_Y=G[V_Y\cup R]$.
  Так как $V_X\cap V_Y=\emptyset$, имеем $|V(G_X)|<|V|$ и $|V(G_Y)|<|V|$.

  \begin{Statement}{}{}
    Любое $XR$-разделяющее множество в $G_X$ является $XY$-разделяющим в $G$.
  \end{Statement}
  \begin{proof}
    Пусть $S$ — $XR$-разделяющее в $G_X$, и в $G\setminus S$ есть путь $p$ из $X$ в $Y$.
    Поскольку $R$ разделяет $X$ и $Y$ в $G$, путь $p$ проходит через некоторую $r\in R$.
    Начальный отрезок $p$ от $X$ до первой $r$ целиком лежит в $G_X$ и не пересекает $S$
    — противоречие.
  \end{proof}

  Любое $XR$-разделяющее в $G_X$ имеет размер $\ge k$.
  По ИП: $k$ непересекающихся $XR$-путей в $G_X$ и $k$ непересекающихся $YR$-путей в $G_Y$.
  Склеиваем через $R$ (возможно, так как $V_X\cap V_Y=\emptyset$) и получаем $k$ путей $X\to Y$.

  \medskip
  \noindent\textbf{Случай 2.}
  Любое $XY$-разделяющее $R$ с $|R|=k$ поглощает $X$ или $Y$.

  \noindent\textit{(а) $X\cup Y \neq V$.}
  Возьмём $u \in V\setminus(X\cup Y)$.

  \begin{Statement}{}{}
    В $G\setminus\{u\}$ нет $XY$-разделяющего множества размера $k-1$.
  \end{Statement}
  \begin{proof}
    Если такое $S$ существует, то $S' = S\cup\{u\}$ — $XY$-разделяющее в $G$
    с $|S'|=k$ и $X\setminus S'\neq\emptyset$, $Y\setminus S'\neq\emptyset$
    (т.\,к. $u\notin X\cup Y$) — противоречие случаю~2.
  \end{proof}

  В $G\setminus\{u\}$ все разделяющие имеют размер $\ge k$, и $|V(G\setminus\{u\})|<|V|$,
  поэтому по ИП есть $k$ путей в $G\setminus\{u\}$, а значит и в $G$.

  \noindent\textit{(б) $X\cup Y = V$.}
  Каждая вершина $X\cap Y$ входит в любое разделяющее.
  Граф $H = G[V\setminus(X\cap Y)]$ двудольный с долями $X'=X\setminus Y$
  и $Y'=Y\setminus X$.
  По теореме Кёнига:
  \[
    k = |\text{мин. покрытие } H| + |X\cap Y|
      = |\text{макс. паросочетание } H| + |X\cap Y| \]
   \[
      = \text{число непересекающихся } XY\text{-путей.}
  \]
\end{proof}

\subsection{Теорема Менгера}

\begin{Theorem}{Менгера}{}
  Максимальное число попарно непересекающихся по \emph{внутренним} вершинам
  путей между несмежными $v$ и $u$ равно минимальному размеру
  $vu$-разделяющего множества $S \subset V\setminus\{v,u\}$.
\end{Theorem}

\begin{proof}
  Пусть минимальный разделитель имеет размер $k$.
  Заменим $v$ на $k$ клонов $v_1,\ldots,v_k$ (с теми же соседями) и $u$ на $u_1,\ldots,u_k$.
  Положим $X=\{v_1,\ldots,v_k\}$, $Y=\{u_1,\ldots,u_k\}$.
  Теорема Геринга даёт $k$ непересекающихся $XY$-путей,
  что соответствует $k$ внутренне непересекающимся $vu$-путям в $G$.
\end{proof}

\begin{figure}[H]
\centering
\begin{tikzpicture}[
  >=Stealth, font=\small,
  v/.style={circle, draw, fill=white, minimum size=8mm, inner sep=0pt},
  vs/.style={circle, draw=blue!70!black, fill=blue!20, minimum size=8mm, inner sep=0pt},
  vt/.style={circle, draw=red!70!black,  fill=red!20,  minimum size=8mm, inner sep=0pt}]

\node[vs] (s) at (0, 0)   {$s$};
\node[v]  (a) at (2, 1.2) {$a$};
\node[v]  (b) at (2,-1.2) {$b$};
\node[v]  (c) at (4, 1.2) {$c$};
\node[v]  (d) at (4,-1.2) {$d$};
\node[vt] (t) at (6, 0)   {$t$};

\draw[gray!40] (a)--(b) (c)--(d) (a)--(d) (b)--(c);
\draw[blue!70!black, very thick] (s)--(a)--(c)--(t);
\draw[red!70!black,  very thick, dashed] (s)--(b)--(d)--(t);

\node[blue!70!black] at (3,  1.9) {путь 1};
\node[red!70!black]  at (3, -1.9) {путь 2};
\end{tikzpicture}
\caption{Теорема Менгера ($k=2$): два попарно непересекающихся по внутренним вершинам
  пути $s\!\to\!t$ соответствуют минимальному разделяющему множеству размера 2.}
\end{figure}
